
\documentclass{article}
\usepackage[utf8]{inputenc}
\usepackage{geometry}
\usepackage{hyperref}
\usepackage{enumitem}

\geometry{a4paper, margin=25mm}

\title{Project Report: Computer Science Portfolio}
\author{Karandeep Singh}
\date{\today}

\begin{document}

\maketitle

\section{Introduction}
This report outlines the design, implementation, and content of my personal computer science portfolio, developed as part of the requirements for the B201 Computer Science Lab module [cite: 99]. As a Computer Science undergraduate at Gisma University of Applied Sciences, I built this portfolio to showcase my technical skills, educational background, and academic projects to prospective employers [cite: 104, 110, 111]. The project is hosted live on GitHub Pages, with all source documentation maintained within my GitHub repository [cite: 110].

\section{Portfolio Structure and Design}
The portfolio is designed as a modern, interactive, and responsive single-page website [cite: 104]. It features sections for Home, About, Technical Skills, Academic Projects, an Interactive Terminal, and Contact information [cite: 104]. The website uses a dark-themed, glass-morphism aesthetic to provide a professional and high-performance visual experience [cite: 104]. The corresponding GitHub repository is organized to house the website's source code, my professional CV in LaTeX format, and documentation related to my Python, SQL, and Data Structures and Algorithms exercises [cite: 104, 110]. This structure allows hiring managers to easily navigate both my web presence and the technical depth of my work [cite: 110].

\section{Tools and Technologies}
The portfolio was developed using several key technologies [cite: 110]. The frontend is built with HTML5 and CSS3, incorporating interactive elements managed by JavaScript [cite: 104]. I utilized GitHub Pages for hosting, which enables efficient deployment of static sites [cite: 104, 110]. My CV, provided in the accompanying \textit{karanCV.pdf}, was typeset using LaTeX, reflecting a commitment to high-quality documentation [cite: 104, 110]. Version control is managed entirely through Git and GitHub, ensuring that all project iterations are documented [cite: 104, 110].

\section{Reflections and Future Improvements}
Developing this portfolio provided practical experience in web development and professional presentation [cite: 110]. A core strength is the interactive terminal feature, which demonstrates my technical creativity and ability to create engaging user experiences [cite: 104]. While the current site is fully functional, it relies on static content [cite: 110]. Future improvements will focus on integrating more dynamic backend functionalities and expanding the showcase of complex technical projects to further demonstrate my evolving skill set [cite: 110].

\section{Submission Links}
In accordance with the requirements detailed in the \textit{B201 Assessment Brief (4).pdf}, the following links provide access to my work:
\begin{itemize}
    \item \textbf{Portfolio Website:} \url{https://karandeep6595.github.io/} [cite: 104]
    \item \textbf{GitHub Repository:} \url{https://github.com/karandeep6595/} [cite: 104]
\end{itemize}

\end{document}
